\chapter{{\fontsize{14pt}{21pt}\selectfont\MakeUppercase{Внедрение и эксплуатация программного комплекса}}}
\label{cha:impl}

\section{Развёртывание программного комплекса}

\subsection{Состав поставки}

В состав поставки программного комплекса входят исходный код монорепозитория (каталоги \texttt{backend}, \texttt{web}, \texttt{mobile}), файл конфигурации контейнеризации \texttt{docker-compose.yml}, файлы зависимостей (\texttt{requirements.txt}, \texttt{package.json}, \texttt{pubspec.yaml}), миграции схемы базы данных (\texttt{Alembic}), скрипт начального заполнения демонстрационными данными (\texttt{scripts.seed\_demo}), а также сопровождающая документация в каталоге \texttt{docs} (архитектура, runbook развёртывания, стратегия тестирования).\cite{docker_compose_docs,fastapi_docs,flutter_docs,react_docs}

\subsection{Требования к аппаратному и программному обеспечению}

\textbf{Сервер развёртывания (минимальные требования):} процессор с архитектурой x86\_64, не менее 2 виртуальных ядер; оперативная память не менее 4 Гбайт; дисковое пространство не менее 20 Гбайт для образов контейнеров, базы данных и медиафайлов; ОС Linux или Windows с поддержкой Docker Engine и Docker Compose; сетевая связность для доступа клиентов по протоколам HTTP/HTTPS.\cite{docker_compose_docs,postgresql_docs}

\textbf{Сервер (рекомендуемые требования):} 4 ядра и 8 Гбайт ОЗУ при одновременной работе десятков пользователей и хранении медиаконтента; отдельный том или сетевое хранилище для каталога медиафайлов, смонтированного в контейнер серверной части.

\textbf{Рабочее место администратора веб-панели:} современный браузер с поддержкой ECMAScript 2015+, разрешение экрана не ниже 1366×768 пикселей.\cite{react_docs,typescript_docs}

\textbf{Мобильное устройство слушателя:} смартфон или планшет под управлением Android или iOS; для сборки отладочного APK --- установленный SDK платформы и среда Flutter.\cite{flutter_docs}

\subsection{Пошаговая процедура развёртывания}

Развёртывание в контейнерном контуре выполняется в следующей последовательности:

\begin{enumerate}
  \item подготовить среду исполнения: установить Docker Engine и Docker Compose, обеспечить доступ к портам \texttt{127.0.0.1:5433} (PostgreSQL), \texttt{127.0.0.1:8000} (серверная часть), \texttt{127.0.0.1:8081} (веб-панель);
  \item клонировать репозиторий и перейти в корень проекта;
  \item задать переменные окружения (минимально --- пароль СУБД \texttt{POSTGRES\_PASSWORD}, секрет приложения \texttt{LMS\_SECRET\_KEY}, режим \texttt{LMS\_ENVIRONMENT}, список разрешённых источников CORS \texttt{LMS\_CORS\_ORIGINS}) по образцу \texttt{.env.example};
  \item выполнить сборку и запуск сервисов командой \texttt{docker compose up --build};
  \item дождаться состояния готовности контейнера СУБД (healthcheck \texttt{pg\_isready}), после чего серверная часть выполнит миграции \texttt{alembic upgrade head} при старте контейнера;
  \item при необходимости загрузить демонстрационные данные профилем \texttt{tools}: \texttt{docker compose run --rm --profile tools seed};
  \item проверить доступность API эндпоинтов состояния: \texttt{GET /health} и \texttt{GET /health/db}, выполнить пробный вход в веб-панель по адресу веб-сервиса.\cite{docker_compose_docs,fastapi_docs,postgresql_docs}
\end{enumerate}

На рисунке~\ref{fig:ch3-deployment} приведена схема локального развёртывания: контейнеризированные серверная часть, веб-панель и СУБД; мобильный клиент подключается к API по сети хоста.

\begin{figure}[H]
\centering
\includegraphics[width=\textwidth,height=0.72\textheight,keepaspectratio]{backend}
\caption{Схема локального развёртывания программного комплекса}
\label{fig:ch3-deployment}
\end{figure}

\subsection{Развёртывание мобильного клиента}

Мобильное приложение не входит в состав Docker Compose и собирается средствами Flutter. Для отладки на эмуляторе Android базовый адрес API задаётся как \texttt{http://10.0.2.2:8000/api/v1}; для физического устройства --- IP-адрес хоста в локальной сети. Переопределение выполняется параметром \texttt{--dart-define=API\_BASE=...}. После сборки APK приложение устанавливается на устройство слушателя; проверка выполняется авторизацией под учётной записью слушателя и запросом к защищённым маршрутам.\cite{flutter_docs,fastapi_docs}

\subsection{Типовые проблемы развёртывания}

В таблице~\ref{tab:ch3-deploy-troubleshoot} приведены типовые симптомы, причины и действия администратора.

\begin{table}[H]
\centering
\caption{Типовые проблемы развёртывания и способы их устранения}
\label{tab:ch3-deploy-troubleshoot}
{\setstretch{1}\fontsize{12pt}{14pt}\selectfont
\setlength{\tabcolsep}{4pt}
\setlength{\extrarowheight}{2pt}
\renewcommand{\arraystretch}{1}
\begin{tabularx}{\textwidth}{|>{\RaggedRight\arraybackslash}p{38mm}|>{\RaggedRight\arraybackslash}p{42mm}|>{\RaggedRight\arraybackslash}X|}
\hline
\textbf{Симптом} & \textbf{Вероятная причина} & \textbf{Действие} \\
\hline
Контейнер серверной части перезапускается циклически & ошибка подключения к СУБД или миграций & просмотреть журнал \texttt{docker compose logs backend}; проверить \texttt{LMS\_DATABASE\_URL}, доступность \texttt{db}, состояние тома PostgreSQL \\
\hline
\texttt{401 Unauthorized} при запросах с заголовком организации & не выбран арендатор или недействительный JWT & выполнить вход, проверить заголовок выбора организации и срок действия токена \\
\hline
Веб-панель не загружает данные & неверный базовый URL API или CORS & проверить сборочный аргумент \texttt{VITE\_API\_BASE} и переменную \texttt{LMS\_CORS\_ORIGINS} \\
\hline
Мобильный клиент не достигает API & неверный \texttt{API\_BASE} или файрвол & задать корректный IP хоста; открыть порт 8000 для LAN \\
\hline
\end{tabularx}
}
\end{table}

Подробные инструкции для администратора и пользователя мобильного приложения приведены в приложении~Г.

\section{Опытная эксплуатация и тестирование}

\subsection{План тестирования}

План испытаний включает функциональное тестирование REST API и сценариев ролей, интеграционное тестирование связки серверная часть --- СУБД --- веб-клиент, нагрузочное тестирование критических эндпоинтов (выборочно), юзабилити-оценку административной панели и мобильного клиента, а также проверку базовых мер безопасности (изоляция данных арендаторов, RBAC, заголовки CORS).\cite{ulikov_testing_2024,moodle_lms,iso29119,kaner_testing_en}

\subsection{Покрытие критического пути}

В таблице~\ref{tab:ch3-critical-path} зафиксировано соответствие шагов критического пути (см. главу~1) статусам прохождения в ходе опытной эксплуатации. Фактические даты и номера сборок --- % TODO(Лусинэ): указать номер сборки и дату прогона при подготовке к защите.

\begin{table}[H]
\centering
\caption{Покрытие критического пути пользователей в ходе опытной эксплуатации}
\label{tab:ch3-critical-path}
{\setstretch{1}\fontsize{12pt}{14pt}\selectfont
\setlength{\tabcolsep}{4pt}
\begin{tabularx}{\textwidth}{|>{\centering\arraybackslash}p{12mm}|>{\RaggedRight\arraybackslash}X|>{\centering\arraybackslash}p{28mm}|}
\hline
\textbf{Шаг} & \textbf{Сценарий} & \textbf{Статус} \\
\hline
1 & Вход преподавателя или администратора организации, выбор организации & пройден \\
\hline
2 & Создание курса, уроков, теста и вопросов & пройден \\
\hline
3 & Назначение курса слушателю или группе & пройден \\
\hline
4 & Прохождение урока и запуск адаптивного теста со стороны слушателя & пройден \\
\hline
5 & Завершение попытки, расчёт результата и формирование рекомендаций & пройден \\
\hline
6 & Просмотр аналитики преподавателем или администратором & пройден \\
\hline
\end{tabularx}
}
\end{table}

\subsection{Метрики опытной эксплуатации}

Для объективной оценки стабильности фиксировались: суммарное число прогонов автоматизированных тестов серверной части и web; число выявленных дефектов по категориям (функциональные, регрессионные, дефекты документации); среднее время отклика типовых REST-методов; доля успешно завершённых попыток адаптивного теста без ошибок серверной части. Для субъективной оценки удобства могут использоваться опросные шкалы SUS; значения --- % TODO(Лусинэ): заполнить после проведения опроса экспертов или пользователей.

\subsection{Сопоставление прогнозируемых и фактических характеристик}

В таблице~\ref{tab:ch3-metrics-compare} приведено сопоставление целевых показателей этапа проектирования и результатов опытной эксплуатации. Численные значения времени отклика и пропускной способности --- % TODO(Лусинэ): вставить результаты замеров (например, из CI или локального профилирования).

\begin{table}[H]
\centering
\caption{Сопоставление прогнозируемых и фактических характеристик}
\label{tab:ch3-metrics-compare}
{\setstretch{1}\fontsize{12pt}{14pt}\selectfont
\setlength{\tabcolsep}{4pt}
\begin{tabularx}{\textwidth}{|>{\RaggedRight\arraybackslash}p{52mm}|>{\centering\arraybackslash}p{38mm}|>{\centering\arraybackslash}p{38mm}|}
\hline
\textbf{Показатель} & \textbf{Прогноз} & \textbf{Факт (опытная эксплуатация)} \\
\hline
Доступность API после развёртывания & без блокирующих ошибок на критическом пути & подтверждено \\
\hline
Среднее время отклика типового GET-запроса & $< 500$ мс при локальной сети & см. замеры \\
\hline
Изоляция данных арендаторов & отсутствие пересечений выборок по \texttt{tenant\_id} & подтверждено тестами \\
\hline
\end{tabularx}
}
\end{table}

Методика и полный перечень тест-кейсов приведены в приложении~Б.

\section{Анализ информационной безопасности}

\subsection{Модель угроз}

Для корпоративной системы управления обучением характерны следующие классы угроз: несанкционированный доступ к учётным записям и токенам сессии; попытки межарендаторского доступа к данным (нарушение изоляции по \texttt{tenant\_id}); инъекционные атаки на уровне СУБД при ошибках формирования запросов; несанкционированное повышение привилегий (обход RBAC); перехват канала без TLS; утечки персональных данных пользователей.\cite{yojji_multitenant_lms_2025,moodle_multitenancy_2024}

\subsection{Реализованные защитные меры}

В разрабатываемом комплексе применяются: разделение ролей (\texttt{learner}, \texttt{teacher}, \texttt{org\_admin}, \texttt{system\_admin}); выдача JWT доступа и обновления с ограниченным временем жизни; хранение хэшей паролей вместо открытого текста; фильтрация запросов к данным по идентификатору арендатора; валидация входных схем через механизм моделей запросов; журнал аудита ключевых действий; ограничение источников CORS в эксплуатационной конфигурации.\cite{fastapi_docs,pydantic_docs,django_rest_framework}

\subsection{Соответствие ФЗ-152 и практикам OWASP}

Обработка персональных данных пользователей (адрес электронной почты, ФИО, результаты обучения) должна выполняться с учётом требований законодательства о персональных данных: минимизация состава данных, разграничение доступа, возможность блокировки учётной записи.\cite{metodichka_vkr_ikt_2021}

Рекомендуется размещение системы в контуре с использованием HTTPS/TLS на границе, регулярное обновление зависимостей и контроль конфигурации по контрольным спискам OWASP Top 10 и OWASP ASVS.\cite{owasp_top10_2021,fz152}

\subsection{Ограничения и направления усиления}

В текущей версии не реализованы отказоустойчивый кластер СУБД, аппаратный модуль безопасности для хранения секретов, шифрование медиаконтента на уровне тома и полноценный Web Application Firewall. Перспективные меры: внедрение OIDC/SAML для единого входа, ротация ключей подписи JWT, политика резервного копирования и восстановления.

\section{Экономическое обоснование разработки}

\subsection{Трудозатраты по фазам}

В таблице~\ref{tab:ch3-effort} приведена оценка трудозатрат по фазам жизненного цикла в человеко-часах. Числовые значения --- ориентировочные для учебного проекта; уточнение --- % TODO(Лусинэ): согласовать с фактическим журналом работ или отчётом по практике.

\begin{table}[H]
\centering
\caption{Оценка трудозатрат по фазам работ}
\label{tab:ch3-effort}
{\setstretch{1}\fontsize{12pt}{14pt}\selectfont
\setlength{\tabcolsep}{4pt}
\begin{tabularx}{\textwidth}{|>{\RaggedRight\arraybackslash}p{62mm}|>{\centering\arraybackslash}p{28mm}|>{\centering\arraybackslash}p{36mm}|}
\hline
\textbf{Фаза} & \textbf{Длительность (нед.)} & \textbf{Трудозатраты (чел.-ч)} \\
\hline
Анализ предметной области и аналогов & 4 & 120 \\
\hline
Проектирование архитектуры и модели данных & 5 & 160 \\
\hline
Реализация серверной части и API & 8 & 320 \\
\hline
Реализация веб-панели и мобильного клиента & 8 & 280 \\
\hline
Тестирование, документация, подготовка демонстрации & 4 & 100 \\
\hline
\textbf{Итого} & \textbf{29} & \textbf{980} \\
\hline
\end{tabularx}
}
\end{table}

\subsection{Себестоимость разработки}

Себестоимость складывается из фонда оплаты труда (ставка разработчика --- % TODO(Лусинэ): указать ставку ₽/ч для Москвы 2026 г.), накладных расходов организации (если применимо), аренды вычислительных ресурсов для эксплуатации стенда и лицензий ПО. Компоненты с открытым исходным кодом (FastAPI, PostgreSQL, React, Flutter, Docker) не образуют лицензионных отчислений при соблюдении условий соответствующих лицензий.\cite{fastapi_docs,postgresql_docs,flutter_docs,react_docs}

\subsection{Альтернативные затраты на коммерческую LMS}

Для организации численностью порядка 200 сотрудников альтернативой является закупка коммерческой корпоративной LMS или облачной подписки. Ориентировочные годовые затраты на подписку зависят от числа активных пользователей и набора модулей и подлежат уточнению по прайс-листам поставщиков --- % TODO(Лусинэ): вставить 2--3 ориентира с официальных сайтов и датой обращения.

\subsection{Срок окупаемости}

Пусть $C_{\mathrm{разр}}$ --- совокупные затраты на разработку и внедрение собственного комплекса, $C_{\mathrm{альт}}$ --- годовые затраты на коммерческую LMS при сопоставимой функциональности, $E$ --- ожидаемая годовая экономия за счёт автоматизации процессов обучения и администрирования (в денежном выражении). Тогда окупаемость в годах может быть оценена соотношением
\begin{equation}
T_{\mathrm{ок}} = \frac{C_{\mathrm{разр}}}{C_{\mathrm{альт}} + E}.
\label{eq:payback}
\end{equation}

Подстановка численных значений --- % TODO(Лусинэ): выполнить после уточнения $C_{\mathrm{альт}}$ и $E$.

\section{Маркетинговая стратегия и план развития продукта}

\subsection{Целевые сегменты}

Целевыми сегментами являются средний корпоративный сектор (100--1000 сотрудников), образовательные организации с собственным контуром обучения, корпоративные академии и учебные центры крупных компаний.\cite{educause_adaptive_2016,docebo_adaptive_2026}

\subsection{Каналы продвижения}

Рекомендуется сочетание технической демонстрации на профильных мероприятиях, публикации кейсов внедрения и партнёрских программ с HR-tech интеграторами.\cite{mobile_learning_solutions_2024}

\subsection{Модель монетизации и лицензирование}

На перспективу возможна модель SaaS-подписки по числу активных пользователей и числу организаций-арендаторов либо лицензирование коробочной поставки для развёртывания на инфраструктуре заказчика с платной технической поддержкой. Открытый исходный код учебной части может распространяться на условиях лицензии MIT; расширения для корпоративных интеграций целесообразно выделять в отдельный проприетарный модуль --- это соответствует практике двойного лицензирования в индустрии ПО.\cite{metodichka_vkr_ikt_2021}

\subsection{План развития на 12 месяцев}

В горизонте 12 месяцев планируются: интеграция единого входа по протоколам OIDC/SAML; расширение статистических моделей адаптивности до параметрических моделей IRT (двух- и трёхпараметрических); развитие мобильного клиента под iOS в составе единого кода Flutter; элементы геймификации; подключение внешних видеоконференций; поэтапная приоритизация изменений по практикам гибкой разработки.\cite{discovery_intelligent_adaptive,flutter_docs,sutherland2016scrum}

\section{Выводы по главе}

В третьей главе рассмотрены вопросы развёртывания программного комплекса в контейнерном окружении и работы мобильного клиента в связке с REST API. Описаны минимальные и рекомендуемые требования к серверным и клиентским средствам, приведена процедура запуска и загрузки демонстрационных данных, таблица типовых сбоев и действий администратора. Сформулирован план испытаний и показатели опытной эксплуатации с опорой на критический путь пользователей. Выполнен анализ информационной безопасности с учётом мультиарендной модели данных и ролевого доступа. Дано экономическое обоснование с формулой оценки срока окупаемости и обозначены маркетинговые сегменты и направления развития продукта. Результаты главы подтверждают готовность комплекса к демонстрации и дальнейшему развитию в составе корпоративной образовательной среды.
